\section{Avatar Forcing}\label{sec:method}
In this work, we present Avatar Forcing, which generates a video of a real-time interactive head avatar, conditioned on the avatar audio and the multimodal signals of the user. We provide an overview of our framework in~\cref{fig:architecture1}.

In \cref{sec:df_motion_sample}, we present our framework for achieving real-time interactive head avatar generation based on causal diffusion forcing in the motion latent space. This consists of two key steps: encoding multimodal user and avatar signals, and causal inference of avatar motion.
In \cref{sec:dpo}, we introduce a preference optimization method for the interactive motion generation, which enhances expressive interaction without the need for additional human labels.

%%%%%%%%%%%%%%%%%%%%%%%%%%%%%%%%%%%%%%%%%%%%%%%%%%%%%%%%%%%%%%%%%%%%%%%%%%%%%%%%%%%%%%%%%%%%%%%%
\begin{figure*}[t]
    \centering
    \vspace{-0.1in}
    \includegraphics[width=\linewidth]{figure/architecture1-4.pdf}
    \vspace*{-8mm}
    \caption{\textbf{Overall architecture of Avatar Forcing}. We encode the use motion and audio, as well as avatar audio into a unified condition by Dual Motion Encoder. Causal Motion Generator infer the motion latent block of the avatar, which are then decoded into an avatar video.} \label{fig:architecture1}
    \vspace*{-2mm}
\end{figure*}
%%%%%%%%%%%%%%%%%%%%%%%%%%%%%%%%%%%%%%%%%%%%%%%%%%%%%%%%%%%%%%%%%%%%%%%%%%%%%%%%%%%%%%%%%%%%%%%%

%%%%%%%%%%%%%%%%%%%%%%%%%%%%%%%%%%%%%%%%%%%%%%%%%%%%%%%%%%%%%%%%%%%%%%%%%%%%%%%%%%%%%%%%%%%%%%%%
\begin{figure}[t]
    \centering
    \vspace{-0.05in}
    \includegraphics[width=0.49\textwidth]{figure/architecture2-2.pdf}
    \vspace*{-6mm}
    \caption{\textbf{Architecture of} $v_{\theta}$. The look-ahead causal attention mask enables a smooth transition across the blocks.}
    \label{fig:architecture2}
\vspace{-0.2in}
\end{figure}
%%%%%%%%%%%%%%%%%%%%%%%%%%%%%%%%%%%%%%%%%%%%%%%%%%%%%%%%%%%%%%%%%%%%%%%%%%%%%%%%%%%%%%%%%%%%%%%%

\subsection{Diffusion Forcing for Real-Time Interaction} \label{sec:df_motion_sample}
\paragraph{Motion Latent Encoding} 
For motion encoding, we employ the motion latent auto-encoder from \citet{float}. The latent auto-encoder maps an input image $S\in\real^{3 \times\! H \!\times\! W}$ to a latent $z \in \real^d$ whose identity and motion are decomposable:
\begin{equation}
    z = z_{S} + \mathbf{m}_S \in \real^d,
\end{equation}
where $z_{S}\in\real^d$ and $\mathbf{m}_S\in\real^d$ are the identity and motion latents, respectively. The identity latent $z_S$ encodes identity representation of the avatar image $S$ (e.g., appearance), while the motion latent $\mathbf{m}_S$ encodes rich verbal and non-verbal features (e.g., facial expression, head motion). We use this latent representation to capture holistic head motion and fine-grained facial expression, which are crucial for realistic head avatar generation.
We provide more details on the auto-encoder in Appendix~\ref{sec:apx_architecture}.


\paragraph{Interactive Motion Generation}
In Avatar Forcing, the motion latents are generated by conditioning on multimodal user signals and avatar audio. This can be formulated as an autoregressive model as follows:
\begin{equation}
    \vspace*{-1mm}
    p_{\theta}( \mathbf{m}^{1:N}) = \prod_{i=1}^{N} p_{\theta}(\mathbf{m}^i | \mathbf{m}^{<i}, \mathbf{c}^{\leq i}), \label{eq:task}
    \vspace*{-1mm}
\end{equation}
where each motion latent $\mathbf{m}^i$ is predicted from past motion latents and the condition triplet $\mathbf{c}^{i} = (\mathbf{a}_u^{i}, \mathbf{m}_u^i, \mathbf{a}^i)$ consisting of user audio $\mathbf{a}_u^{i}$, user motion $\mathbf{m}_u^i$, and avatar audio $\mathbf{a}^i$.


Based on this formulation, we introduce a diffusion forcing-based causal motion generator operating in the motion latent space, which is modeled using a vector field model $v_\theta$. As illustrated in~\cref{fig:architecture2}, the model $v_\theta$ comprises two main components: Dual Motion Encoder and Causal DFoT Motion Generation.

The goal of the Dual Motion Encoder is to capture the bidirectional relationship between multimodal user signals and avatar audio, and to encode them into a unified condition. As illustrated in~\cref{fig:architecture2}(a), the encoder first takes the user signals ($\mathbf{m}^{i}_u$ and $\mathbf{a}^{i}_u$) and aligns them through a cross-attention layer, which captures the holistic user motion. This representation is then integrated with the avatar audio using another cross-attention layer, which learns the causal relation between the user and the avatar, producing a unified user-avatar-condition. 
In \cref{sec:experiments:ablation}, we empirically validate the importance of using user motion $\mathbf{m}^{i}_u$ for generating an interactive avatar.

For the causal motion generator, we adopt the diffusion forcing transformer (DFoT)~\citep{historydiffusion} with a blockwise causal structure~\citep{causvid, ar_wo_vq}. The latent frames are divided into blocks to capture local bidirectional dependencies within each block while maintaining causal dependencies across different blocks. 
For each block, we assign a shared noise timestep to all frames and apply an attention mask that prevents the current block from attending to any future blocks.

In \cref{fig:blockwise_causal_dit}, we compare our motion generator architecture with the standard bidirectional DiT architecture used in INFP~\citep{infp}. Unlike INFP, which requires the full temporal context, our diffusion forcing allows stepwise motion generation under causal constraints and user-avatar interaction with low latency. 

%%%%%%%%%%%%%%%%%%%%%%%%%%%%%%%%%%%%%%%%%%%%%%%%%%%%%%%%%%%%%%%%%%%%%%%%%%%%%%%%%%%%%%%%%%%%%%%%
\begin{figure}[t]
    \centering
    \vspace{-0.05in}
    \includegraphics[width=0.48\textwidth]{figure/cache4.pdf}
    \vspace*{-7mm}
    \caption{\textbf{Architectural comparison between bidirectional and causal structure}. (a) Bidirectional DiT used in INFP~\citep{infp} requires access to the entire temporal window for motion generation. (b) Our blockwise causal DFoT predicts the next block without using future context and supports KV caching.
} \label{fig:blockwise_causal_dit}
\vspace{-0.15in}
\end{figure}
%%%%%%%%%%%%%%%%%%%%%%%%%%%%%%%%%%%%%%%%%%%%%%%%%%%%%%%%%%%%%%%%%%%%%%%%%%%%%%%%%%%%%%%%%%%%%%%%

However, we observe that the strict causal mask in the blockwise causal attention fails to ensure a smooth temporal transition across blocks.
To address this issue, we introduce \emph{look-ahead} in the causal mask, which allows each block to attend to a limited number of future frames while preserving overall causality.
We define this blockwise look-ahead causal mask $M$ as follows (illustrated in~\cref{fig:architecture2}(b), right):
\begin{align}
    M_{i,j} = 1 ~~\text{if}~~ \lfloor j / B \rfloor \leq \lfloor i / B \rfloor + l ~~\text{else}~~ 0, \label{eq:lookahead_mask}
\end{align}
where $i$ and $j$ are the frame indices, $B$ denotes the block size, and $l$ is the look-ahead frame size. With look-ahead, we effectively mitigate severe per-frame motion jittering observed in the simple blockwise causal structure, which we provide video examples in the supplementary materials.


Based on the blockwise causal structure, the motion generator takes a noisy latent block as input, which is concatenated with the avatar motion latent $\mathbf{m}_S$ serving as the reference motion.
We use a cross-attention layer to condition on the unified condition from the Dual Motion Encoder. 
Additionally, we employ a sliding-window attention mask of size $2l$ along the time axis for temporal smooth conditioning. We provide the more details on $v_{\theta}$ in Appendix~\ref{sec:apx_architecture}.

%%%%%%%%%%%%%%%%%%%%%%%%%%%%%%%%%%%%%%%%%%%%%%%%%%%%%%%%%%%%%%%%%%%%%%%%%%%%%%%%%%%%%%%%%%%%%%%%
\begin{figure}[t]
\vspace{-0.15in}
\centering
\begin{minipage}{1.0\linewidth}
\begin{algorithm}[H]
\small 
\caption{\label{algo:sampling} Motion inference with KV caching} 
\begin{algorithmic}[1]
\Require ODE timesteps $\{t_n\}^T_{n=0}$, motion generator $v_{\theta}$, video length $N$, block size $B$, max cache size $M$, user inputs $(\mathbf{a}_{u},\mathbf{m}_{u})$, avatar audio $\mathbf{a}^i$, latent-to-frame decoder $\texttt{Dec}$, and id latent $z_S$. 
\State \textbf{Divide} the frames into $L = \lceil N / B \rceil$ blocks
\State \textbf{Initialize} $\mathbf{KV}, \mathbf{cKV}$ $\leftarrow [~], [~]$ $\rhd$ Frame \& condition caches
\For{$i = 1$ to $L$}
    \State \textbf{Sample} Noise block $\mathbf{m}_{t_{0}}^{i} \sim \mathcal{N}(0, I)$, \hfill $\rhd$ $\mathbf{m}_{t_{0}}^{i}\in\real^{B \times d}$
    \State \textbf{Acquire} \highlightline{User inputs ($\mathbf{a}^i_{u}$, $\mathbf{m}_{u}^i$)
    and avatar audio $\mathbf{a}^i$}.
    \State \textbf{Set} $\mathbf{c}^i$ $\leftarrow$ $(\mathbf{a}^i_{u}, \mathbf{m}_{u}^i, \mathbf{a}^i)$ \hfill $\rhd$ Condition triplet
    \For{$j = 0$ to $T$}
        \State \textbf{Solve ODE:} $\mathbf{m}_{t_{j+1}}^{i} \leftarrow v_{\theta}(\mathbf{m}_{t_j}^{i}, t_j; \mathbf{c}^i, \mathbf{KV}, \mathbf{cKV})$
    \EndFor
    \State \textbf{Decode} \& \textbf{Return}~~ \highlightline{$\mathbf{x}_1^i$ $\leftarrow$ $\texttt{Dec}(z_S, \mathbf{m}_1^i)$~$\in$~$\real^{B \times 3 \times H \times W}$}
    \State \textbf{Update caches}~~ $\mathbf{kv}_i$, $\mathbf{ckv}_i$ $\leftarrow$ $v_{\theta}(z_{1}^{i}, 1; \mathbf{c}^i, \mathbf{KV}, \mathbf{cKV})$    
    \If{$|\mathbf{KV}|$~$=$~$|\mathbf{cKV}|$~$=$~$M$}
    \State $\mathbf{KV}{\texttt{.pop(0)}}$~and~  $\mathbf{cKV}{\texttt{.pop(0)}}$
    \EndIf
    \State $\mathbf{KV}${\texttt{.append($\mathbf{kv}_i$)}}~and~ $\mathbf{cKV}${\texttt{.append($\mathbf{ckv}_i$)}}
\EndFor
\end{algorithmic}
\end{algorithm}
\end{minipage}
\vspace{-0.3in}
\end{figure}
%%%%%%%%%%%%%%%%%%%%%%%%%%%%%%%%%%%%%%%%%%%%%%%%%%%%%%%%%%%%%%%%%%%%%%%%%%%%%%%%%%%%%%%%%%%%%%%%

\paragraph{Training and Inference}
To train the vector field model $v_{\theta}$, we formulate the diffusion forcing training objective in Eq.~\eqref{eq:train_df} into a motion latent generation objective as follows:
\begin{equation}
	\loss_{DF}(\theta) = \mathbb{E}_{n, t_n, \mathbf{m}^n_{t_n}}\! \| v_{\theta}(\mathbf{m}^n_{t_n}, t_n, \mathbf{c}^n) \!-\! (\mathbf{m}_1^n \!-\! \mathbf{m}_0^n)\|,
    \label{eq:train_ours}
\end{equation}
where $n$~$\in$~$[1,N]$ denotes the frame index, $t_n$~$\in$~$[0,1]$ is the per-frame noise timestep, $\mathbf{c}^n = (\mathbf{a}_u^{i}, \mathbf{m}_u^i, \mathbf{a}^i)$ is the user-avatar condition triplet, $\mathbf{m}_{t_n}^n$~$\in$~$\real^d$ is the noisy motion latent. For simplicity, we omit the reference motion condition $\mathbf{m}_S$ in Eq.~\eqref{eq:train_ours}.

With the trained model, we can generate the interactive avatar motion given user inputs and avatar audio in an autoregressive rollout. We provide the pseudo code for motion sampling in \cref{algo:sampling}, which adopts rolling KV cache in a blockwise manner~\citep{causvid, self_forcing}. Further details on inference are provided in Appendix~\ref{sec:apx_exp_detail}.

\subsection{Enhancing Motion Expressiveness} \label{sec:dpo}
While our model enables real-time motion generation, we observe that it struggles to produce expressive motions crucial for natural human conversation. In contrast to achieving an accurate lip-sync that has an almost one-to-one match with the avatar audio, appropriate reaction to the user's motion and audio is highly ambiguous, as there is no single \emph{correct} response~\citep{l2l, diffusion_listening, customlistener}. 

Moreover, as shown in \cref{fig:vis_speaker_listener_exp}, existing listening datasets generally exhibit lower motion expressiveness compared to speaking datasets, which leads models to learn passive and non-expressive listening behaviors.
While prior work attempted to address it with personalized motion~\citep{l2l} or external text instructions~\citep{customlistener}, these methods do not address the core challenge of learning the true interactive behavior for natural interaction.

To generate vibrant and expressive reactions, we formulate this as an alignment problem and apply the Reinforcement Learning from Human Feedback (RLHF)~\citep{ouyang2022training} approach. 
The main challenge lies in defining an explicit reward for the avatar’s interactive behavior, which must account for the interplay between avatar audio, user audio, user motion, and the generated video, which are difficult to evaluate even for humans.

We adapted a reward-free method, Direct Preference Optimization (DPO)~\citep{rafailov2023dpo, wallace2024diffusiondpo}, for fine-tuning our motion generation based on diffusion forcing. Specifically, we construct preference pairs of motion latents $(\mathbf{m}^w, \mathbf{m}^l)$ as follows:
\begin{itemize}[leftmargin=1.5em]
\item \textbf{Preferred sample} $\mathbf{m}^w$: the motion latent from the ground-truth video, exhibiting expressive and contextually appropriate responses.
\item \textbf{Less preferred sample} $\mathbf{m}^l$: the motion latent generated by a separately trained \textit{talking avatar model}~\citep{float}, conditioned solely on the avatar audio.
\end{itemize}
This paired design yields a clear signal for enhancing expressiveness, emphasizing active listening and reactive motion while leaving other aspects, such as lip sync or speech-driven motion, unchanged.

Building on the original DPO objective $\loss_{\text{DPO}}$ (Eq.~\eqref{eq:dpo}) with our constructed pairs $(\mathbf{m}^w, \mathbf{m}^l)$, we fine-tune our diffusion forcing model $v_{\theta}$ with the following objective:

\begin{equation}
    \loss_{ft}(\theta) = \loss_{DF}(\theta) + \lambda\loss_{DPO}(\theta), \label{eq:total_objective}
\end{equation}
where $\lambda$ is a balancing coefficient. As a result, our model achieves efficient preference alignment without requiring a dedicated reward model, which we validate in \cref{sec:experiments:ablation}. We provide the detailed formulation of $\loss_{\text{DPO}}$ in Appendix~\ref{sec:apx_dpo}.

%%%%%%%%%%%%%%%%%%%%%%%%%%%%%%%%%%%%%%%%%%%%%%%%%%%%%%%%%%%%%%%%%%%%%%%%%%%%%%%%%%%%%%%%%%%%%%%%
\begin{figure}[t]
\centering
    \vspace{-0.05in}
    \includegraphics[width=\linewidth]{figure/exp_l2_var_exp-pasted.pdf}
    \vspace{-0.3in}
    \caption{\textbf{Variance visualization} of the L2-norm of 3DMM expressions~\citep{spectre} for the speaker and listener on ViCo~\citep{rlhg} dataset. Higher variance indicates higher expressiveness.}\label{fig:vis_speaker_listener_exp}
\vspace{-0.15in}
\end{figure}
%%%%%%%%%%%%%%%%%%%%%%%%%%%%%%%%%%%%%%%%%%%%%%%%%%%%%%%%%%%%%%%%%%%%%%%%%%%%%%%%%%%%%%%%%%%%%%%%